\chapter{Loss of Bone Density (Osteopenia/Osteoporosis)}

\section*{Definition}
Progressive reduction in bone mineral density, leading to weakened skeletal structure and increased fracture risk, most commonly resulting from oestrogen deficiency after menopause.

\section*{Pathophysiology}
Oestrogen inhibits osteoclast activity through receptor activator of nuclear factor-kB ligand (RANKL) regulation. Oestrogen binds to oestrogen receptors on osteoclasts, promoting apoptosis and reducing bone resorption. The rapid decline in oestrogen during menopause removes this inhibitory signal, uncoupling bone remodelling: osteoclast activity increases while osteoblast activity does not keep pace, leading to net bone loss. Trabecular bone (vertebrae, distal radius) is affected more rapidly than cortical bone. Peak bone loss occurs in the first 3-5 years after menopause.

\section*{Etiology}
\begin{itemize}
  \item Menopausal oestrogen deficiency (primary cause)
  \item Advanced age
  \item Low peak bone mass (genetic, nutritional, exercise-related)
  \item Calcium and vitamin D deficiency
  \item Glucocorticoid use (corticosteroid-induced osteoporosis)
  \item Smoking
  \item Excessive alcohol consumption
  \item Low body weight (BMI < 19)
  \item Family history of osteoporosis or hip fracture
  \item Secondary causes (hyperthyroidism, hyperparathyroidism, malabsorption, rheumatoid arthritis)
\end{itemize}

\section*{Indications for Evaluation}
Seek care if:
\begin{itemize}
  \item Severe or worsening symptoms
  \item Bone density screening (DXA scan) is recommended for all women aged 65 and older, or younger with risk factors; seek advice after any low-trauma fracture
  \item Accompanied by other concerning signs
\end{itemize}

\section*{Related Symptoms}
\begin{itemize}
  \item Back pain
  \item Loss of height
  \item Joint pain
  \item Fractures
  \item Muscle weakness
\end{itemize}

\textbf{Note:} Always consider serious underlying conditions when evaluating persistent loss of bone density.

\section*{Self-Care / Home Management}
\begin{itemize}
  \item Rest and avoid activities that may aggravate the condition.
  \item Maintain adequate hydration and a balanced diet.
  \item Over-the-counter remedies may provide symptomatic relief (consult a pharmacist or doctor).
  \item Monitor symptoms and seek professional advice if they persist or worsen.
  \item Avoid self-medicating with prescription medications without medical guidance.
\end{itemize}

\section*{Diagnostic Approach}
\begin{itemize}
  \item A thorough history and physical examination are the first steps in evaluation.
  \item Laboratory tests (blood work, urinalysis) may be ordered based on clinical suspicion.
  \item Imaging studies (X-ray, ultrasound, CT, MRI) may be indicated when structural pathology is suspected.
  \item Specialist referral is arranged when the diagnosis remains uncertain or requires specific expertise.
\end{itemize}

\section*{Treatment / Management Overview}
\begin{itemize}
  \item Treatment targets the underlying cause identified through diagnostic evaluation.
  \item Supportive care and symptom management are provided while awaiting definitive therapy.
  \item Medications, lifestyle modifications, and physical therapies may be prescribed as appropriate.
  \item Follow-up ensures treatment effectiveness and allows timely adjustments to the care plan.
\end{itemize}

\clearpage